\documentclass[12pt,a4paper]{article}

% ---------- Мова та шрифти ----------
\usepackage{fontspec}
\usepackage{polyglossia}
\setmainlanguage{ukrainian}
\setotherlanguage{english}
\setmainfont{DejaVu Serif}
\setsansfont{DejaVu Sans}
\setmonofont{DejaVu Sans Mono}

% ---------- Компонування ----------
\usepackage[a4paper,margin=2.2cm]{geometry}
\usepackage{graphicx}
\usepackage{float}
\usepackage{caption}
\usepackage{xcolor}
\usepackage{hyperref}
\usepackage{enumitem}
\usepackage{titlesec}
\usepackage{longtable}
\usepackage{booktabs}
\usepackage{array}
\usepackage{amssymb}
\usepackage{fancyvrb}

\hypersetup{colorlinks=true, linkcolor=blue!50!black, urlcolor=blue!50!black,
    pdftitle={UML-діаграми проєкту --- Хмарне сховище}, pdfauthor={Проніна Марина}}

\titleformat{\section}{\normalfont\Large\bfseries}{\thesection.}{0.6em}{}
\titleformat{\subsection}{\normalfont\large\bfseries}{\thesubsection}{0.6em}{}
\setlength{\parindent}{0pt}
\setlength{\parskip}{0.6em}

\newcommand{\fig}[4]{%
    \begin{figure}[H]
        \centering
        \includegraphics[width=#3\linewidth,height=0.85\textheight,keepaspectratio]{#1}
        \caption{#2}
        \label{#4}
    \end{figure}
}

\begin{document}
\begin{titlepage}
\centering\large
\textbf{Київський національний університет імені Тараса Шевченка}\\
Факультет комп'ютерних наук та кібернетики
\vspace*{\stretch{1}}

\textbf{Звіт}\\
з дисципліни “Інформаційні Технології”\\
до Практичного завдання -- модуль 1\\
Варіант 1.2\\
“Етап 1. UML-специфікація”
\vspace*{\stretch{1}}

\begin{flushright}
Студентки 4 курсу\\
групи МІ-41\\
спеціальності 122 "Комп’ютерні науки"\\
Проніної Марини
\end{flushright}
\vspace*{\stretch{1.5}}

\textbf{Київ \the\year}
\end{titlepage}

\begin{abstract}
\noindent
У документі наведено комплект UML-діаграм для проєкту \textbf{«Хмарне сховище»} --- програмної системи типу 2 (клієнт віддаленої файлової папки, спрощений аналог Google Drive) за варіантом \textbf{12}: перегляд вмісту файлів \texttt{.xml} (як текст) і \texttt{.png} (як зображення), сортування за датою модифікації за зростанням/спаданням та фільтрація «усі файли / .java / .png». Система складається з сервера на Python (FastAPI + SQLite), десктоп-клієнта (PySide6/Qt) та веб-клієнта (Vue 3). Діаграми описують систему такою, якою її реально реалізовано, і впорядковані за життєвим циклом проєктування: від аналізу вимог до деталізованого проєктування підсистем. Усі діаграми побудовано в нотації Mermaid; вихідні тексти --- у \texttt{docs/\allowbreak{}diagrams-uml/\allowbreak{}src/\allowbreak{}}, зображення --- у \texttt{docs/\allowbreak{}diagrams-uml/\allowbreak{}img/\allowbreak{}}.
\end{abstract}

\tableofcontents
\newpage


% =====================================================================
\section{Аналіз вимог}
% =====================================================================

Номер студентського квитка закінчується на \textbf{12}: перша цифра (1) задає типи файлів для перегляду вмісту, друга (2) --- номер у списку ОПЕРАЦІЙ. Для системи типу 2: ТИПИ №1 --- \texttt{.xml}, \texttt{.png}; ОПЕРАЦІЇ №2 --- «сортування за датою і часом модифікації (зростання і спадання); фільтр для відображення всіх файлів або лише файлів типу .java, .png». Реалізація:

{\small
\begin{longtable}{@{}p{5cm}p{11cm}@{}}
\caption{Вимоги варіанта 12}\\
\toprule
\textbf{Вимога} & \textbf{Реалізація} \\
\midrule
\endhead
Перегляд вмісту при кліку (ТИПИ №1) & \texttt{.xml} --- як текст (моноширинний шрифт), \texttt{.png} --- як зображення \\ \midrule
Сортування (ОПЕРАЦІЇ №2) & за датою та часом модифікації, за зростанням $\blacktriangle$ і спаданням $\blacktriangledown$ \\ \midrule
Фільтрація (ОПЕРАЦІЇ №2) & «Усі файли» / «Лише .java» / «Лише .png» \\ \midrule
Бонус (+5) & drag-n-drop вивантаження у веб- і десктоп-клієнті \\ \midrule
\bottomrule
\end{longtable}}

Загальні функціональні вимоги до системи типу 2:

\begin{itemize}[nosep]
    \item реєстрація та автентифікація на віддаленому сервері (логін/пароль, JWT-токен);
    \item персональний простір кожного користувача --- інші користувачі не бачать і не можуть змінити чужі файли;
    \item список файлів з атрибутами: ім'я, дата створення, дата модифікації, хто вивантажив, хто редагував;
    \item показ/приховування будь-якого стовпця (дати, автори, розмір), крім імені (стовпець «Ім'я» приховати неможливо);
    \item вивантаження, завантаження та видалення файлів; повторне вивантаження файлу з тим самим іменем --- це редагування (оновлюються вміст, дата модифікації та «Редагував»);
    \item синхронізація обраної локальної папки з віддаленою (у обидва боки, з виявленням конфліктів).
\end{itemize}

Архітектурно система --- клієнт-серверна: REST API на FastAPI зберігає метадані у SQLite, а вміст файлів --- у каталозі \texttt{data/\allowbreak{}blobs/\allowbreak{}} через абстракцію \texttt{FileStorage}. Обидва клієнти працюють з одним і тим самим API. Логіка варіанта (сортування, фільтр, типи перегляду) та алгоритм синхронізації винесені в спільний пакет \texttt{cloudfolder.core}, а у веб-клієнті продубльовані в модулях \texttt{variant.js} і \texttt{sync.js} з ідентичною поведінкою, що підтверджується однаковими тестовими наборами.


% =====================================================================
\section{Архітектурне проєктування (високий рівень)}
% =====================================================================

Діаграми побудовано в порядку життєвого циклу ПЗ: після аналізу вимог спершу описано архітектуру системи загалом (прецеденти, класи високого рівня, розгортання), а далі --- деталізоване проєктування окремих модулів (класи VOPC, послідовності, стани, компоненти, діяльність).

\subsection{Діаграма прецедентів}

Діаграма (рис. 1) показує основного актора --- \emph{Користувача} --- та допоміжного актора \emph{Таймер / ОС}, який ініціює автоматичну синхронізацію при змінах у локальній папці. Усі операції з файлами \emph{включають} («include») прецедент «Автентифікуватися». Прецеденти варіанта --- сортування за датою модифікації та фільтр «усі / .java / .png» --- \emph{розширюють} («extend») перегляд списку; перегляд вмісту розширюється двома альтернативами відображення (.xml / .png). Drag-n-drop --- розширення вивантаження.

\emph{Примітка щодо нотації:} Mermaid не має окремого типу діаграми прецедентів, тому її побудовано засобами \texttt{flowchart}: прецеденти --- овали, межа системи --- \texttt{subgraph}, стереотипи підписано на зв'язках.

\fig{img/01-usecase.png}{Діаграма прецедентів системи «Хмарне сховище»}{0.95}{fig:1}

\subsection{Діаграма класів (загальна структура)}

Рис. 2 відображає доменну модель і шари сервера: сутності \texttt{User} та \texttt{FileRecord} (файл має SHA-256 вмісту, що потрібен для синхронізації), сервіси \texttt{AuthService} і \texttt{FileService}, репозиторії над \texttt{Database} (SQLite) та інтерфейс сховища \texttt{FileStorage} з реалізацією \texttt{LocalFileStorage}. Модуль \texttt{variant} інкапсулює правила варіанта і використовується як сервером (параметри \texttt{sort}, \texttt{order}, \texttt{type} у \texttt{GET /\allowbreak{}api/\allowbreak{}files}), так і десктоп-клієнтом.

\fig{img/02-class-highlevel.png}{Діаграма класів високого рівня}{0.95}{fig:2}

\subsection{Діаграма розгортання}

Рис. 3 описує фізичну топологію: на ПК користувача виконуються десктоп-клієнт (Python 3.11) та браузер із SPA; сервер --- Docker-контейнер на платформі Render, у якому процес uvicorn одночасно обслуговує REST API (HTTPS) та роздає статичну збірку веб-клієнта (\texttt{web/\allowbreak{}dist}); поруч --- файл бази SQLite і каталог з вмістом файлів.

\fig{img/14-deployment.png}{Діаграма розгортання}{0.95}{fig:3}


% =====================================================================
\section{Деталізоване проєктування}
% =====================================================================

\subsection{VOPC-діаграма для прецедента варіанта «Сортувати та фільтрувати список»}

Як прецедент для VOPC-діаграми (\emph{View Of Participating Classes}) обрано операцію варіанта --- сортування за датою модифікації та фільтрацію «усі / .java / .png». Класи розподілено за стереотипами аналізу:

\begin{itemize}[nosep]
    \item «boundary» --- \texttt{MainWindow} (комбобокси фільтра й сортування, клік по заголовку таблиці) та \texttt{FilesView} у веб-клієнті;
    \item «control» --- \texttt{FilesViewModel} (стан сортування/фільтра, без залежності від Qt) і модуль \texttt{variant};
    \item «entity» --- \texttt{FileRecord}; допоміжний перелік \texttt{SortOrder} (none $\to$ $\blacktriangle$ $\to$ $\blacktriangledown$ $\to$ none);
    \item «interface» --- \texttt{ApiClientProtocol}, завдяки якому view-модель тестується із заглушкою замість мережі.
\end{itemize}

\fig{img/03-class-vopc-sort-filter.png}{VOPC-діаграма прецедента «Сортувати та фільтрувати список»}{0.95}{fig:4}

\subsection{VOPC-діаграма для прецедента «Синхронізувати локальну папку»}

Синхронізація --- найскладніший сценарій. \texttt{SyncEngine} отримує локальний маніфест (\texttt{scan\_folder}), віддалений список і базовий стан із файлу \texttt{.cloudfolder-sync.json}, будує план чистою функцією \texttt{plan\_sync} і виконує дії через \texttt{ApiClient}. Подію зміни папки надає \texttt{QFileSystemWatcher}.

\fig{img/04-class-vopc-sync.png}{VOPC-діаграма прецедента «Синхронізувати локальну папку»}{0.95}{fig:5}

\subsection{Діаграми взаємодії (послідовності)}

Рис. 6--9 показують взаємодію об'єктів у часі для автентифікації, вивантаження (з розгалуженням «новий файл / редагування наявного»), операції варіанта та синхронізації.

\fig{img/07-seq-login.png}{Діаграма послідовності: автентифікація}{0.95}{fig:6}

\fig{img/08-seq-upload.png}{Діаграма послідовності: вивантаження файлу}{0.95}{fig:7}

\fig{img/09-seq-sort-filter.png}{Діаграма послідовності: фільтр «Лише .java» та сортування за датою модифікації}{0.95}{fig:8}

\fig{img/10-seq-sync.png}{Діаграма послідовності: синхронізація папки}{0.95}{fig:9}

\subsection{Діаграми станів}

Рис. 10 --- життєвий цикл окремого файлу з погляду синхронізації; рис. 11 --- стани клієнтської сесії (вкладений стан «Автентифіковано» містить режими роботи з файлами).

\fig{img/11-state-file-sync.png}{Діаграма станів файлу під час синхронізації}{0.95}{fig:10}

\fig{img/12-state-session.png}{Діаграма станів клієнтської сесії}{0.95}{fig:11}

\subsection{Діаграма компонентів}

Рис. 12 показує програмні компоненти: два клієнти, спільну бібліотеку \texttt{cloudfolder.core} та сервер. Клієнти взаємодіють із сервером лише через REST-інтерфейс \texttt{/\allowbreak{}api/\allowbreak{}*} (JSON + Bearer JWT).

\fig{img/13-component.png}{Діаграма компонентів}{0.95}{fig:12}

\subsection{Діаграми діяльності}

Рис. 13 описує цикл синхронізації: тристороннє порівняння \emph{local / remote / base} дає одну з п'яти дій (upload, download, delete\_remote, delete\_local, conflict). Конфлікт (файл змінено з обох боків) розв'язується без втрати даних: локальна версія зберігається як «ім'я (конфлікт).ext» і вивантажується, а віддалена записується під основним іменем. Помилка однієї дії не зупиняє цикл.

\fig{img/05-activity-sync.png}{Діаграма діяльності: цикл синхронізації}{0.95}{fig:13}

Рис. 14 описує операцію варіанта: обробку зміни фільтра або порядку сортування.

\fig{img/06-activity-sort-filter.png}{Діаграма діяльності: сортування та фільтрація}{0.95}{fig:14}


% =====================================================================
\section{Відповідність вимогам етапу}
% =====================================================================

{\small
\begin{longtable}{@{}p{5cm}p{3cm}p{3cm}p{5cm}@{}}
\caption{Кількість діаграм за типами}\\
\toprule
\textbf{Тип діаграми} & \textbf{Вимога} & \textbf{Виконано} & \textbf{Файли} \\
\midrule
\endhead
Прецедентів & 1+ & 1 & 01-usecase \\ \midrule
Класів (з VOPC) & 2+, $\geq$1 VOPC & 3 (2 VOPC) & 02, 03, 04 \\ \midrule
Діяльності & 1+ & 2 & 05, 06 \\ \midrule
Взаємодії (послідовності) & 2+ & 4 & 07--10 \\ \midrule
Станів & 1+ & 2 & 11, 12 \\ \midrule
Компонентів & 1+ & 1 & 13 \\ \midrule
Розгортання & 1+ & 1 & 14 \\ \midrule
\bottomrule
\end{longtable}}


% =====================================================================
\section{Висновки}
% =====================================================================

Побудовано 14 UML-діаграм семи типів, що повністю покривають вимоги етапу 1. Діаграми узгоджені з фактичним кодом: назви класів, методів і REST-маршрутів збігаються з реалізацією у \texttt{cloudfolder/\allowbreak{}} та \texttt{web/\allowbreak{}src/\allowbreak{}}. Винесення правил варіанта та алгоритму синхронізації в окремі чисті модулі (\texttt{variant}, \texttt{sync}) спростило як проєктування (VOPC-діаграми мають чітко відокремлені «control»-класи), так і модульне тестування.


\end{document}
